\documentclass[10pt,a4paper]{article}

% --- Layout & typography (ultra-tight, 3 columns) ---
\usepackage[margin=0mm]{geometry}
\usepackage{multicol}
\setlength{\columnsep}{0pt}
\setlength{\columnseprule}{0pt}
\pagestyle{empty}
\setlength{\parskip}{0pt} % no vertical glue between paragraphs

\usepackage{titlesec,needspace}
\titlespacing*{\section}{0pt}{2pt plus 1pt}{1pt}
\titlespacing*{\subsection}{0pt}{1pt}{1pt}

\usepackage{enumitem}
\setlist{leftmargin=*, itemsep=1pt, topsep=2pt, partopsep=0pt}
\usepackage{booktabs}
\usepackage[hidelinks]{hyperref}
\hypersetup{colorlinks=true,linkcolor=black,urlcolor=black,citecolor=black,
  pdftitle={CS1101S Cheatsheet}, pdfauthor={} }

% --- Code highlighting (minted) ---
\usepackage{minted,xcolor}
\setminted{
  fontsize=\scriptsize,
  breaklines,
  frame=single,   % frame on to visually separate sections
  framesep=1pt,   % tight inner padding
  linenos=false,
  tabsize=2
}
% --- Narrower fonts + global 9pt body ---
\usepackage{newtxtext}                 % narrower serif body
\usepackage[scaled=0.92]{inconsolata}  % compact monospace for code
\IfFileExists{microtype.sty}{\usepackage[stretch=10,shrink=10]{microtype}}{}
\AtBeginDocument{\fontsize{6}{10.5}\selectfont}  % 9pt body (10.5pt leading)

% Zero vertical gap around minted blocks
\linespread{0.80}
% Remove vertical list glue around minted blocks
\usepackage{etoolbox}
\AtBeginEnvironment{minted}{%
  \setlength{\parskip}{0pt}%
  \setlength{\topsep}{0pt}%
  \setlength{\partopsep}{0pt}%
  \setlength{\parsep}{0pt}%
  \vspace{-2pt}%
}

\AtEndEnvironment{minted}{\vspace{-2pt}}

% --- Document ---
\begin{document}\footnotesize

\begin{multicols*}{3}
\raggedcolumns

% ===================== 1) Variables =====================
\begin{minted}{javascript}
//T is a subtype of S, T<:S
  //variables of type T can be used on code written for type S
//widening conversion
  //primitives: type with smaller range to a type with larger range
  //reference types: more specific type to more general type
//narrowing conversion
  //java does not allow narrowing unless explicitly casted
//byte <: short <: int <: long <: float <: double
//char <: int <: long <: float <: double
\end{minted}
% ===================== 2) Functions =====================
\begin{minted}{javascript}
return_type function_name(type1 param1, type2 param2){...}
//if does not return anything, use void
\end{minted}
% ===================== 3) Encapsulation =====================
\begin{minted}{javascript}
//Encapsulation: keeping objects data and code that works on data together, hiding internal details from other code
  //other code interacts with class through methods
  //should not directly mess with objects fields
//objects of a given class called instance of a class
//when created, class fields assigned default values
  //whole numbers assigned 0, real numbers 0.0, boolean false, char \u0000, object reference null
//aliasing: 2 reference variables can share the same value
//accessing null fields and methods cause NullPointerException
\end{minted}
% ===================== 4) Private and Public =====================
\begin{minted}{javascript}
//private only access within class, public can access/modify outside class
\end{minted}
% ===================== 5) Constructors =====================
\begin{minted}{javascript}
//Constructors invoked automatically when instantiated, MUST have same name as class, NO return type
  //runs when new is used
//this keyword
  //reference variable that refers back to the calling object itself
//if no constructor given, default constructor will automatically be added at CT
  Circle(){}
\end{minted}
% ===================== 6) Accessors and Mutators =====================
\begin{minted}{javascript}
//tell, dont ask principle
  //tell an object what you want done, instead of asking for its data and doing the work yourself
\end{minted}
% ===================== 7) Class Fields =====================
\begin{minted}{javascript}
class Math {
  public static final double PI = 3.14159;
}
//static keyword: associate field with class in java
//only 1 shared copy available, can be accessed by any instance
//can also be accessed with className
//final keyword: does not allow the value of the field to be changed
  //final instance field must be assigned value by the end constructor else compile error
//use public static final for constant values
\end{minted}
% ===================== 8) Class Methods =====================
\begin{minted}{javascript}
//static method: does not depend on any particular object instance
class Circle{
  public static int getCircleCount() {...}
}
//CANNOT use instance fields in static method
  //can use class fields in instance method
public static final void main(String[] args) {...}
\end{minted}
% ===================== 9) Composition =====================
\begin{minted}{javascript}
//Composition (has-a relationship)
  //just storing an object in class
\end{minted}
% ===================== 10) Heap and Stack =====================
\begin{minted}{javascript}
//Stack: local variables and stack frames
  //instance and class fields are NOT variables
  //call frames are created when method is invoked and removed when method completes
  //LIFO (last element inserted is first to be removed)
//Heap: stores all objects
  //whenever new keyword is used



//during instance method call, JVM will create a stack frame
  //it contains *this* reference, method arguements and local variables within method
//if class method called, no *this* reference
\end{minted}
% ===================== 11) Inheritence =====================
\begin{minted}{javascript}
class ColoredCircle extends Circle {
  private Color color;
  public ColoredCircle(Point center, double radius, Color color) {
    super(center, radius);  // call the parent's constructor
    this.color = color;
  }} //ColouredCircle <: Circle
//all public fields and methods in parent will be in child, including public class fields/methods
//NO private fields
//if constructor does not invoke the superclass constructor, compiler will insert default constructor
  //if superclass has constructor and not called, CT error
Circle c = new ColoredCircle(p, 0, blue); //widening type conversion
ColoredCircle cc = c; //compile error, compiler treats c as Circle, dk that its RTT is ColouredCircle
ColoredCircle cc = (ColoredCircle) c; // narrowing type conversion
  //compiler treats c as if it has CTT ColouredCircle
  //if wrong, ClassCastException at RT
//ALL classes inherit from the Object class
  //boolean Object::equals(Object)
  //String Object::toString()
//if child extends Parent, you can call inherited static method using the child className too
\end{minted}
% ===================== 12) Overriding =====================
\begin{minted}{javascript}
//method signature: name, type, order and number of parameters
  class C {
    A foo(B1 x, B2 y) {...}} // C::foo(B1, B2)
//method descriptor: method signature + return type
  //A C::foo(B1, B2)
//overriding: must have same method descriptor as instance method in parent class
  //return type may be a subtype of the og return type
  //parameters MUST be exactly the same
  //dc about parameter names
  //CANNOT override class method
class ColouredCircle() {
  ...
  @Override
  public String toString() {...}
}
@Override
  //hint to compiler that method is intended to override method in parent class
  //if type or overriding not possible, compiler will lyk by throwing error
//toString method, String Object::toString()
  //converts a reference object to String object
  //called automatically during string concatenation using the + operator
  String s = "Circle c is" + c; //Circle c is Circle@1ce92674
  //can override toString method, so it returns different string, will still be called automatically
\end{minted}
% ===================== 13) Overloading =====================
\begin{minted}{javascript}
//Overloading: class has access (defined or inherited) to 2 or more methods with SAME name but different method signatures
  //either order, number or type of parameters different, dc about names
  //if only difference is names or no difference at all between 2 methods, compile error
  //return type can be different, CANNOT differ only by return type
class Circle {
  public boolean contains(Point p) {...}
  public boolean contains(Point p, double y){...}
}
//constructor can also be overloaded
  class Circle {
    public Circle(Point c, double r) {...}
    public Circle() { this(new Point(0, 0), 1); }
  }  //one constructor can call another constructor in the same class
//CAN overload static class methods
  //can also overload inherited public class methods
\end{minted}
% ===================== 14) Polymorphism =====================
\begin{minted}{javascript}
//Polymorphism: change how existing code behaves without changing any existing code
  //achieved through method overriding
//overriding equals method
class Circle {
  ...
  @Override
  public boolean equals(Object obj) {
    if (obj instanceof Circle) {
      Circle circle = (Circle) obj;
      return (circle.c.equals(this.c) && circle.r == this.r);
    }
    return false;
  }} //== DOES NOT call equals method in java
//obj instanceof operator
  //check if runtime type of obj is a subtype of Circle, obj<:Circle
\end{minted}
% ===================== 15) Method invocation =====================
\begin{minted}{javascript}
//dynamic binding
boolean same = curr.equals(obj)
//method signature being called (decided at CT)
  //1. determine the CTT of curr
  //2. compiler search for all methods in curr (defined and inherited) that can be correctly invoked with arguement type
    //also searches widening references (subtype to supertype) or primitives (int to double)
  //3. if more than 1 method, choose most specific one based on arguements
  //4. if fails to determine a single most specific method, throw compile error
  //5. compiler stores method descriptor in bytecode
  //CT chooses the most specific applicable overload
//class that provides the implementation (decided at RT)
  //1. method descriptor from CT is retrieved
  //2. determine the RTT of curr, let it be T
  //3. starting from T, select nearest method implementation with matching method descriptor in class hierachy until it reaches root Object
//ONLY applies to instance methods, class methods DO NOT support dynamic binding
  //CT part is the same, then java will invoke the method defined in CTT of target, ignores the RTT
\end{minted}
% ===================== 16) Liskov substitution principle (LSP) =====================
\begin{minted}{javascript}
//Liskov substitution principle (LSP): subclass should not break expectations set by superclass
  //subclass should be able to pass all test cases of superclass, else not substitutable and LSP violated
//even though S<:T, code meant for T will not always work on S
  //S might override some methods from T, and output from those methods might be different even though return type is the same
//final keyword: prevent class or method from being inherited
  final class Circle {...}
  //or can allow inheritence but prevent a specific method from being overriden
  class Circle { public final boolean contains(Point p) {...} }
//final keyword can also be used when declaring local variables to prevent it from being reassigned within method body
\end{minted}
% ===================== 17) Abstract class =====================
\begin{minted}{javascript}
//Abstract class (is-a relationship)
  //class that cannot be instantiated, usually represents an abstract concept
  //if instantiated, will produce an error
abstract class GraphicObject {
  public int x;
  public int y;
  GraphicObject(int x, int y) {
      this.x = x;
      this.y = y;
  }  // A constructor - child classes will call this
  void moveTo(int newX, int newY) {
      this.x = newX;
      this.y = newY;
      System.out.println("Moved to coordinates: " + x + "," + y);
  }  // Concrete method: Every shape moves the same way
  abstract String draw(int x);   // Abstract methods
  abstract void resize(); //no method body
}
class Circle extends GraphicObject {
  Circle(int x, int y) {
      super(x, y); // Calls the GraphicObject constructor
  }
  @Override
  String draw(int level) {
      System.out.println("Drawing a circle at " + x + "," + y);
  } //must have same param and return type
  @Override
  void resize() {
      System.out.println("Resizing the circle.");
  }
} //Circle is-a GraphicObject relationship
//subclass either implement ALL abstract methods or be declared abstract as well
  //class that inherits must implement ALL abstract methods in parent and grandparent
//NOT all methods have to be abstract, can have no abstract method
//class with 1 abstract method MUST be declared abstract
//can contain class methods, call with className
//CANNOT have private abstract and final abstract methods, compile error
//concrete class: non-abstract class
\end{minted}
% ===================== 18) Interface =====================
\begin{minted}{javascript}
//Interface (can-do relationship)
  //models what an entity can do, name usually ends with -able suffix
interface GetAreable {
  double getArea();
} //all methods in interface declared public abstract by default
class Flat extends RealEstate implements GetAreable {
  private int numOfRooms;
  @Override
  public double getArea() {...}
}
//concrete class MUST override ALL abstract methods from interface and provide implementation to each
//abstract class dont need, but child that inherits abstract class needs
class SmartPhone implements Camera, Phone, WebBrowser {...}
//can also extend from 1 or more interfaces
  //interface cannot inherit from class
//interface acts as a supertype
  //if class C implements interface I, C<:I
interface I { ... }
class A { ... }
class B implements I { ... }
A a;
B b;
(I) b; //will compile since B <: I
(I) a; //will also compile although wrong
//will compile as there is a possibility that a actually has a RTT of its subclass which implements from I
  //will fail during RT but will compile
class AI extends A implements I{ ... }
A a = new AI(); // AI extends A and implements I
I i = (I) a;    // this would be safe
//if it is final class A, java knows that it definitely does not have a subclass that implements I, compile error
  //ALL non-final classes can be cast to an interface and it will compile
//impure interfaces: more methods are added to interface, all existing classes will NOT compile unless updated
interface GetAreable {
  double getArea();
  default double getAreaInSquareFeet() { //can include default method
    return getArea() * 10.7639;
  }
} //classes will automatically inherit these methods unless they choose to override them
//can have static methods, call with interface name, no override
//can have private methods, only callable within interface
//ALL variables in interface public static final only, constants
//NO constructors
\end{minted}
% ===================== 19) Wrapper class =====================
\begin{minted}{javascript}
//Wrapper class: class version of a primitive type
//java provides a wrapper for each of its primitive types
  //Integer for int, Double for double and so on
//reference types
Integer b = Integer.valueOf(5);
Integer x = new Integer(5);  //use is discouraged
int j = b.intValue(); //use this to access the value
//auto boxing and unboxing
Integer a = 5; //autoboxing: int -> Integer
int j = i; //unboxing: Integer -> int
//unboxing can throw NullPointerException if object is null
//less efficient that primitive types
Double sum = 0.0;
for (int i = 0; i < Integer.MAX_VALUE; i++) {
  sum += i;
}
//wrapper objects are immutable, everytime the the value is updated, a new Double object is created
//should always use equals method when comparing wrapper objects, override so it compares the values not reference type
  //using == will just compare reference
Double d = 4 //compile error
//cannot simultaneously widen int to Double and autobox double to Double
  //can widen then box or box then widen
\end{minted}
% ===================== 20) Runtime class mismatch =====================
\begin{minted}{javascript}

\end{minted}
% ===================== 21) Variance =====================
\begin{minted}{javascript}
//variance of types
//refers to how the subtype relationship between complex types relates to the subtype relationship between components
//covariant if S<:T implies C(S) <: C(T)
//contravariant if S<:T implies C(T) <: C(s)
//invariant if its not covariant or contravariant
//arrays are covariant
  //S<:T, then S[]<:T[]
Integer[] intArray = new Integer[2] { Integer.valueOf(10), Integer.valueOf(20) };
Object[] objArray;
objArray = intArray;
objArray[0] = "Hello!"; // compiles but RT error
//compiler will allow you to put String into objArray since CT Object and String<:Object but it will have RT of Integer[]
\end{minted}
% ===================== 22) Exceptions =====================
\begin{minted}{javascript}
try { ... }
catch (FileNotFoundException e) { ... }
catch (InputMismatchException e) { ... }
catch (NoSuchElementException e) { ... }
finally { ... }
//exceptions are instances that are a subtype of the Exception class
//try blocks: runs first, stops running once exception is thrown
//catch blocks are checked in order
  //first catch block that has an execution type compatible with the type of the thrown exception (type or subtype) will be selected
  //if ExceptionX <: ExceptionY
  catch(ExceptionY e) { ... } 
  catch(ExceptionX e) { ... }
  //will throw compile error as exception would already be caught on ExceptionY, ExceptionX is redundant
//finally block (optional): runs after all the try and catch blocks, usually for cleanup tasks, ALWAYS RUNS
class Circle {
  ... //have to declare that function throws exception with *throws* keyword
  public Circle(Point c, double r) throws IllegalArgumentException {
    if (r < 0) {
      throw new IllegalArgumentException("radius cannot be negative.");
    } //create Exception object and throw it to caller with *throw* keyword
    ...
  }
} //throw causes the method to return immediately
//unchecked exceptions: caused by programmers errors, subclass of RuntimeException, cause runtime errors
  //RuntimeException is subclass of Exception
//checked exceptions: errors that programmer has no control over, subclass of Exception
class Main {
  static FileReader openFile(String filename) { return new FileReader(filename); }
  public static void main(String[] args) { openFile(); }
} //wont compile cause FileNotFoundException is not handled
//checked exceptions are enforced by compiler
  //if method can throw it, must either declare it with throws or catch it inside the method
  //callers must either catch or declare it too
class Main {
  static FileReader openFile(String filename) {
    try { return new FileReader(filename);} 
    catch (FileNotFoundException e) {
      System.err.println("Unable to open " + filename + " " + e); //calls e.toString()
      return null; //must still return something even though theres an error
    } //if block threw exception, then dont need return anything
  } //since block catches exception, must return something of type FileReader or null
  public static void main(String[] args) { openFile(); }
}
//unchecked exceptions do not need to be declared in throws
  //exception can still happen at runtime but caller is not forced to catch it and no compile error if its ignored
  //code can compile but if theres error will still have a RT error
//Passing the buck: the caller of the method that generates an exception does not need to catch it, can pass it to caller
class Main {
  static FileReader openFile(String filename) throws FileNotFoundException { return new FileReader(filename); }
  public static void main(String[] args) {
    try { openFile(); } 
    catch (FileNotFoundException e) { ... }
  }
}
//creating exceptions: inherit from existing exceptions
class IllegalCircleException extends IllegalArgumentException {
  Point center;
  //overloaded the constructor, if this is called center is null
  IllegalCircleException(String message) { super(message); }
  IllegalCircleException(Point c, String message) { ... }
  @Override
  public String toString() { ... } //this will run when trying to get error string
}
//changing the set of checked exceptions that a method throws is a change to its public specifications
  //may require changing its callers
  //overriden methods are only allowed to throw the same or more specific (subtype) exceptions
  //overriden method can also throw nothing, makes the method easier to call
    //code that handles exception will still compile and run, just that it catches nothing
  //NOT allowed to throw new or broader checked exceptions
//unchecked exceptions, even if superclass method throws unchecked exception
  //in override can throw nothing or a different exception
  //dont have to delcare that you are throwing an unchecked exception with the throws keyword
    //can still throw unchecked exception in code and still cause code to immediately return
//the Error class: when program should terminate when there is no way to recover from error
  //when heap or stack is full
//Exception and Error are subclasses of class Throwable, only objects of this class can be thrown and caught by catch blocks
\end{minted}
% ===================== 23) Generics =====================
\begin{minted}{javascript}
//generic types: class or interface that is parameterised by 1 or more type parameters
//type parameter: what is declared in < >
//type variable: the placeholder for type, declared in definition of generic class or method, S or T in examples
//type arguement: concrete type or type variable passed to a generic type or method when instantiated
//parameterized type: instantiated generic type
class Pair<S,T> { //use single capital letter to name each type variable
  private S first;
  private T second;
  public Pair(S first, T second) {
    this.first = first;
    this.second = second;
  }
  public S getFirst() { return this.first; }
  public T getSecond() { return this.second; }
}
//using a generic type, have to pass in type arguements
  //type arguements can be non-generic type, generic type or another type parameter that has been declared (type variable), ONLY reference types
Pair<String,Integer> foo() { return new Pair<String,Integer>("hello", 4); }
Pair<String,Integer> p = foo();
Integer i = (Integer) p.getFirst(); // compile-time error
  //compiler knows that getFirst is bound to a String, can throw a compile error
class DictEntry<T> extends Pair<String,T> { ... } //pass type parameter of a generic type to another generic type
  //if DictEntry<Integer> it will inherit all accessible instance methods and fields from Pair<String, Integer>
class A {
  public static <T> boolean contains(T[] array, T obj) { ... }
} //generic methods
String[] strArray = new String[] { "hello", "world" };
A.<String>contains(strArray, 123); //compile error
//Bounded type parameters
class A {
  public static <T extends GetAreable> T findLargest(T[] array) {
    for (T curr : array) {
      double area = curr.getArea();
      ...
    } //says that T must be a subtype of GetAreable
    ...
  } //used for both class and interface
} // have to do this as this guarantees that T is able to use the getArea() function, else compile error
//using bounded type parameters for declaring generic classes
class Pair<S extends Comparable<S>,T> implements Comparable<Pair<S,T>> {
  ... //generic interface Comparable<T>
  public Pair(S first, T second) {
    this.first = first;
    this.second = second;
  } //contains method int compareTo(T t) for any concrete class that implements the interface 
  @Override
  public int compareTo(Pair<S,T> s1) {
    return this.first.compareTo(s1.first);
  } //S extends Comparable<S>, self referential bound, comparable to itself
} //implements Comparable<Pair<S,T>>, so compareTo() takes in Pair<S,T> and can be properly overriden
//Pair<String, Integer> is considered a concrete type as it is no longer generic
\end{minted}
% ===================== 24) Type Erasure =====================
\begin{minted}{javascript}
//type erasure: erases type parameters and type arguements during compilation after type checking
//erases the types and replaces it with Object
//cast conversions will be inserted at use sites, when you return something of type T basically
Integer i = new Pair<String,Integer>("hello", 4).getSecond(); //will become
Integer i = (Integer) new Pair("hello", 4).getSecond();
//generated by compiler after type checking, after casting is proven to be correct, 0 chance of exceptions
//if parameter is bounded, it is replaced by the bounds
  //If T extends GetAreable, then T is replaced with GetAreable
class A {
  void foo(Pair<Sring, String> p) { ... }
  void foo(Pair<Integer, Integer> p) { ... }
}
//CANNOT overload when the only difference between the arguements is the type arguements, compile error
  //after type erasure, Pair<...> p will be changed to Pair p, duplicate methods
//CANNOT used class type variables in static field and method
  //static fields/methods share by all instances by type variable belongs to an instance and different instances have different type variables
class D<T> { //static methods must announce its own independent type parameter
  public static <T> T foo(T x) { ... }
} //the type variable T for foo is different from the type variable T for class D
String result = D.<String>foo("Hello World"); //call it normally
//generics and arrays CANNOT mix
  //after type erasure, the array will allow all the generic types to adjust the array
  //will cause heap pollutions where a different parameterized type is placed into the array
  //instantiating a genric array, new Pair<String, Integer>[10] will cause compile error
//generic array declaration, T[] with nothing assigned to it is allowed
//type checking -> type erasure -> dynamic binding
//bridge method: type erasure will affect inheritence and method overriding if they interact with generics
  //bridge methods will be generated to maintain polymorphism
class Box<T> { void get(T t) { ... } }
class StringBox extends Box<String> {
  @Override
  void get(String s) { ... }
} //after type erasure, compiler will see Box::get(Object), different method description, breaking polymorphism
//bridging method generated when
  //class is a subtype of a parameterized type and type erasure changes the signature of any inherited method
class StringBox extends Box {
  void get(Object s) { return get((String) s); }
  void get(String s) { return "hello"; }
}
Box<String> b = new StringBox();
b.get("hi");
//after type erasure, will be get(Object), during CT, compiler searches for closest method signature to get(String) in Box, save get(Object)
//during RT, StringBox will call get(Object), if no bridging method it wont be able to call anything, thus the bridging method
\end{minted}
% ===================== 25) Unchecked Warnings =====================
\begin{minted}{javascript}
ArrayList<Pair<String,Integer>> pairList;
pairList = new ArrayList<Pair<String,Integer>>(); // ok
pairList.add(0, new Pair<Double,Boolean>(3.14, true));  // error
ArrayList<Object> objList = pairList;  // error
//when parameterized, ensures type safety by checking for appropriate types during CT
//generics are invariant, no subtyping relationship between ArrayList<Object> and ArrayList<Pair<...>>
//no possibility of heap pollution
class Seq<T> {
  private T[] array;
//not allowed to assign Object[] to T[] without casting, thus (T[])
  public Seq(int size) {
//SuppressWarnings cannot apply to an assignment, only declaration
    @SuppressWarnings("unchecked") //also write comment to explain how its safe to use this
    T[] a = (T[]) new Object[size];
    this.array = a; //declare local variable then assign to array
  } //tells the compiler everything is fine, no need warning
  public void set(int index, T item) {
    this.array[index] = item;
  } //compiler will check type, if anybody
  public T get(int index) {
    return this.array[index];
  }
  // public T[] getArray() {
  //   return this.array;
  // }
} //remove getArray method as people can put stuff into the array from outside
//unchecked warnings: compiler saying it has done what it can but because of type erasure there could be RT error that it cannot prevent
//T[] is Object[] that was casted to become T[], takes in any reference type outside, which is allowed at RT, heap pollution
  Seq<Integer> si = new Seq<>(1);
  Integer[] ints = si.getArray();
  Object[] o = ints; //allowed, Integer[] <: Object[]
  o[0] = "boom"; //allowed at RT since o actually Object[]
//during type checking compiler will still treat this.array as if it has type of T[]
Seq<Integer> si = new Seq<>(10);
String[] s = si.getArray(); //compile error, compiler sees getArray return Integer[]
//if getArray is removed, only way to place item in array is through set method, which must have type T, array guaranteed to be type T
  //during type checking, if set function called with wrong type T, compile error
//Raw types: generic class without <>
void populateSeq(Seq s) {
  s.set(0, 1234);
  s.set("h","e");
}  //can put any type into the class, compiler will not check, NOT allowed in new code
//only allowed to use when checking object type at RT, obj instanceof Seq
\end{minted}
% ===================== 26) Wildcards =====================
\begin{minted}{javascript}
//NO subtyping relationship in generics, Seq<Circle> NOT a subtype of Seq<Shape>, CT error
//there is subtyping relationship in type variables though
//Upper bound wildcard, Seq<? extends T>
public void copyFrom(Seq<? extends T> src) {
  int len = Math.min(this.array.length, src.array.length);
  for (int i = 0; i < len; i++) {
    this.set(i, src.get(i));
  }
}  //generic type, dk exact type but promise its either T or subclass of T
//S<:T, then A<? extends S> <: A<? extends T> (covariance)
  //A which contains either S or subclass of S is A which contains either T or subclass of T
//For any type S, A<S> <: A<? extends S>
  //A<S> is a subclass of A<S and all its subclasses>
//subtype relation between wildcard types different than normal
  //no longer all code of superclass will also work on subclass
  //more like all subtypes of A<? extends S> is subtype of A<? extends T>
//transitive: A<S> <: A<? extends T>
//if method takes in Seq<? extends T>, get must be assigned to T and set must be null if udk whats inside
//wildcard NOT a type variable, only used in method parameters, foo(<? extends Apple> a)
  //function can return Apple or any of its supertypes
//Lower bounded wildcard, Seq<? super T>
//S<:T, then A<? super T> <: A<? super S> (contravariance)
public void copyTo(Seq<? super T> dest) {
  int len = Math.min(this.array.length, dest.array.length);
  for (int i = 0; i < len; i++) {
      dest.set(i, this.get(i));
  }
} //A that contains T or supertype of T is a subtype of A which contains S or supertype of S
//For any type S, A<S> <: A<? super S>
  //A<S> is a subclass of A<S and all its supertypes>
//transitive: A<T> <: A<? extends S>
//if method takes in Seq<? super S>, get must be assigned to Object and set must be S if udk whats inside
//Producer Extends; Consumer Super (PECS)
  //Producer, ? extends T, can safely get T out
  //Consumer, ? super T, can safely put T into it
//Unbounded wildcards, Seq<?>
void foo(Seq<?> seq) {
  x = seq.get(0);
  seq.set(0, y);
} //method will be flexible to take in a Seq of any type
//restrictive, x MUST be Object and y has to be null as Seq can be any type
//Generics are INVARIANT
Seq<Object> a1 = new Seq<String>(0); 
Seq<Object> a2 = new Seq<Integer>(0); //compile error
Seq<?> a1 = new Seq<String>(0); // Does compile
Seq<?> a2 = new Seq<Integer>(0); // Does compile
//Seq<?> is a sequence of object of some specific but unknown type
//Seq<Object> is a sequence of Object instances with type checking by the compiler
//Seq is a sequence of Object instances without type checking
a instanceof Seq<?> //replace raw types with this
List<?>[] a = new List<?>[10]; //compiles
Comparable<?>[] b = new Comparable<?>[10]; //compiles
//A<?> is a refiable type as no type information is lost during type erasure, thus can create arrays
  //creates an array of A[], can store any A object in it
\end{minted}
% ===================== 27) Type Inference =====================
\begin{minted}{javascript}
Pair<String,Integer> p = new Pair<>();
//java infers that p should be instance of Pair<String, Integer> based on CTT of p
public static <S> boolean contains(Seq<? extends S> seq, S obj) { ... }
A.contains(circleSeq, shape) //can call like non-generic method
//Rules for type inferencing
  //S -> T: S <: T <: Object
  //S -> ? extends T: S <: T <: Object
  //S -> ? super T: T <: S, arguement typing
  //T extends S: T <: S
  //return type T and assignment type S, target typing, T <: S
//ALWAYS choose the lowest bound that everything has in common
//when multiple types satisfy constraints, java choose the most general one that satisfies all bounds
public static <T> boolean contains(T[] array, T obj){}
String[] strArray = new String[] { "hello", "world" }; //String<:T<:Object
A.contains(strArray, 123); // Integer<:T<:Object
//T<:Object satisfies all bounds, will compile but runtime error
public <T extends Circle> T foo(Seq<? extends T> seq)
ColoredCircle c = foo(new Seq<GetAreable>());
//T<:ColouredCircle, T<:Circle, GetAreable<:T<:Object, no solution, compile error
\end{minted}
\end{multicols*}
\end{document}